Two carbon atoms are counted as bonded when their distance is below 2.0\,\AA{} (between the equilibrium bond length of 1.42\,\AA{} and the inner cut-off region 1.95--2.25\,\AA{} of \rebo{}; bonds at the tensile instability are 1.7--1.8\,\AA{} long and ruptured pairs relax to more than 2.5\,\AA{} within one increment). Bond sets are compared between consecutive converged states; every lost bond is recorded with the strain, the two atoms and the midpoint of the bond. Persistence is enforced in post-processing (a bond that reappears is not counted as damage). The damage-localisation measure is $L=1-H/H_\mathrm{max}$, where $H=-\sum_i p_i\ln p_i$ is the Shannon entropy of the distribution $p_i$ of damage-event midpoints over an $8\times8$ grid of the periodic cell and $H_\mathrm{max}=\ln 64$; $L=0$ for uniformly distributed damage and $L=1$ if all events fall in one cell. Load-bearing connectivity is tested by percolation of the bond graph across the periodic $x$ boundary (a doubled-cell union-find). Fracture modes are classified automatically from the number of discrete damage steps, the fraction of bonds lost in the largest single step, the post-peak strain range and the post-peak load retention (\texttt{analysis/metrics.py}).
